\subsection{Datasets and shared taxonomy}

Two public datasets are used (Section~\ref{sec:datasets-bg}). Laboro Tomato
has 643 training and 161 validation images, photographed at
3120$\times$4160 or 3024$\times$4032 pixels in portrait or landscape
orientation. Its six classes combine two fruit sizes (large and small) with
three maturity stages. tomatOD has 277 images, of which 55 form its validation
split. Both datasets are used with their own published train/validation
splits, so no image appears in both training and evaluation.

The class names of both datasets are mapped onto three stages: green,
half-ripened and fully ripened. The conversion uses class names because the
numeric class identifiers are inconsistent between the datasets and between
the splits of tomatOD (Section~\ref{sec:annotation}).

Two preprocessing steps are applied. About 40\% of Laboro images carry an EXIF
orientation tag. The stored pixels are in landscape orientation, but the
annotations assume an upright image. The YOLO data loader ignores the tag, so
the boxes do not align with the fruit. We rotate these images on disk before
conversion to align the pixels and boxes. This correction reduced the
ripe-fruit counting MAE of a six-class baseline detector from 0.93 to 0.82. We
then convert the COCO annotations to the YOLO text format. No other
preprocessing is done offline. Images are resized on the fly to an input size
of 640 pixels: at evaluation and inference the longer side is scaled to 640 and
the image is padded to a 640$\times$640 square, which preserves the aspect
ratio, while during training the standard Ultralytics augmentations (mosaic,
scaling, translation, flipping and HSV colour jitter) are applied at the same
input size.

Three YOLO11n detectors are trained for 60 epochs at 640 pixels, with batch
size 16, seed 0, and no early stopping. One detector is trained on Laboro, one
on tomatOD, and one on both datasets. We refer to the detector trained on both
datasets as the \emph{mixed} model. Each detector is evaluated separately on
both validation splits.

\subsection{Metrics}

Detection quality is reported as COCO-style box mAP50--95 and mAP50
(Section~\ref{sec:metrics-bg}). Counting quality is reported as the mean
absolute error between the predicted and the annotated number of fruit, per
image and per maturity stage,
\begin{equation}
 \mathrm{MAE}_{count} = \frac{1}{3N}\sum_{i=1}^{N}\sum_{s=1}^{3}
 \left|\hat{c}_{i,s}-c_{i,s}\right|.
\end{equation}
Here, $N$ is the number of images and $\hat{c}_{i,s}$ and $c_{i,s}$ are the
predicted and annotated numbers of fruit of stage $s$ in image $i$. Its unit is
fruit per image and stage: a counting MAE of 1.25 means that, on average, the
count for each stage in each image is off by 1.25 fruit. The predicted count is
the number of detections of that stage above the confidence threshold.

For the occlusion experiments, detections are sorted from highest to lowest
confidence and paired with annotated boxes without using the maturity labels.
Each detection and annotation can be paired only once. This greedy matching
separates detection recall from the stage accuracy of the matched detections.
Unless stated otherwise, the confidence threshold is 0.25 and the matching
threshold is IoU 0.5.

\subsection{Colour analysis}
\label{sec:colour}

The colour of a fruit is measured in the CIELAB colour space, which separates
lightness ($L^*$) from two colour-opponent axes. The $a^*$ axis runs from green
(negative values) to red (positive values). It is a dimensionless coordinate
defined by the CIE \cite{cie2018colorimetry}; it has no formal bound, and for
colours that an sRGB camera can record it lies roughly between $-90$ and
$+100$. $a^*$ and the related $a^*/b^*$ ratio are standard colour indices for
tomato ripening in postharvest science \cite{lopezcamelo2004color}, because
the change from chlorophyll green to lycopene red moves a fruit mainly along
this axis.

For each annotated box, we calculate the median $a^*$ in its central half. This
region excludes most of the stem and leaves at the box edges. We convert the
pixels from sRGB to CIELAB with the D65 white point. White balance and
illumination affect the whole image, so we subtract the median $a^*$ of the
entire frame from every box in that frame. The frame median is dominated by
foliage, and the resulting value is called background-corrected $a^*$. For each
dataset, we use grid search to fit two cut points on this axis
(green\,|\,half and half\,|\,fully). The search maximises balanced accuracy
against the labels of that dataset. We then apply the cut points from each
dataset to the boxes from the other dataset. The proportion of labels that
change measures the disagreement between the annotation conventions.
